\scanpage{3}
\begin{center}Structure\end{center}

$4^{\circ})$ La donnée d'une sous-catégorie \inserted{\unsure{str.te|str.t|strict.} pleine} de $M$ (motifs effectifs), stable par $\otimes$, par sommes par \deleted{\unsure{sous et par}} sous-\deleted{\unsure{quo}}objets, telle que \inserted{où $\Lambda(-1)$ soit effectif et que} pour tt $E$, $\exists$ entier $n$ tel que $E(-n)=E\otimes\Lambda(-1)^{\otimes n}$ soit ``effectif''. Une telle structure sur $G_{\deleted{\unsure{C|\mathbb Q}}}$ en définit une autre sur $G_C$, et réciproquement, donc un tel élément de structure peut être considéré comme équivalent à la donnée d'un \deleted{partie de \unsure{$\chi^{+}$|$X^+$} $K(G)$ \unsure{base, à savoir}} \deleted{cône \unsure{$\chi$} $K^+(G)$} cône $K^+(G)$ \deleted{l'anneau \illegible{1} des \unsure{caractères|contenu} de $G_{\unsure{C|\mathbb C}}$ $(\chi^++\chi^+\subset\chi^+)$ \unsure{$\chi$} stable par produit, tel que pour tout} dans \inserted{l'anneau des \unsure{car.} \unsure{rat/$\mathbb{Q}$|virt./$\mathbb Q$}} $K(G)$, contenant $0$, $\deleted{\unsure{1|E}}\varepsilon$ stable par multiplication, et
% slanted
\inserted{engendré par \deleted{\unsure{les}} élém. de la base canonique \deleted{\illegible{1}} (classes de repr. irréd sur $\mathbb{Q}$) qu'il contient.}
À priori, ce semble là un élément de structure faisant intervenir la structure sur $Q$. Il y a lieu de poser d'autres axiomes. \inserted{Par exemple,} on voudra que pour tt \deleted{\unsure{M|H}} $\deleted{\unsure{V}}E\in\operatorname{Ob}M^+$ $E$ soit à graduations \uuline{positives}.
% slanted
\margin{De plus on voudra que la structure envisagée provienne d'une structure de \deleted{\unsure{$=$|même}} type donné sur $G^\circ$, \deleted{\unsure{d'où}} \unsure{par} \deleted{\illegible{1}} \unsure{supposant|supprimant} $G$ \unsure{connexe}.}
\margin{\uline{Axiome} Pour tt \deleted{\unsure{entier|degré} \unsure{$n$}} $i$ il n'y a qu'un nb fini de types de motifs effectifs de poids $i$ \ldots\ \textcircled{\scriptsize ???} En tous cas, on veut que les motifs effectifs de poids $0$ soient tels que $G^\circ$ opère en \unsure{eux} \unsure{trivialem.t}.}

Un lien $G_{m\,C}\xrightarrow{i}G_C$ \deleted{positif} \uline{effectif} si pour tout $E\in\operatorname{Ob}M^+$, la graduation correspondante de $E_C$ est à
% archive sheet number: 2
